\documentclass[11pt]{article}
\usepackage[a4paper,margin=2.1cm]{geometry}
\usepackage{booktabs}
\usepackage{graphicx}
\usepackage{amssymb}
\usepackage{amsmath}
\usepackage{float}
\usepackage{caption}
\usepackage{pdflscape}
\usepackage{xcolor}
\usepackage{framed}
\usepackage{enumitem}
\usepackage{listings}
\usepackage[hidelinks]{hyperref}

% ---- palette ----
\definecolor{dlnavy}{RGB}{18,52,86}
\definecolor{dlaccent}{RGB}{176,42,55}
\definecolor{shadecolor}{RGB}{238,243,248}
\definecolor{codebg}{RGB}{246,246,246}
\definecolor{codekw}{RGB}{18,52,86}
\definecolor{codecom}{RGB}{110,120,130}

\setlist[itemize]{leftmargin=1.4em,itemsep=1.5pt,topsep=2pt}
\setlist[enumerate]{leftmargin=1.6em,itemsep=1.5pt,topsep=2pt}

\lstdefinestyle{rstyle}{
  language=R,
  basicstyle=\ttfamily\scriptsize,
  backgroundcolor=\color{codebg},
  keywordstyle=\color{codekw}\bfseries,
  commentstyle=\color{codecom}\itshape,
  stringstyle=\color{dlaccent},
  breaklines=true,
  showstringspaces=false,
  frame=single,
  rulecolor=\color{gray!40},
  framexleftmargin=3pt,
  xleftmargin=3pt,
  aboveskip=4pt,belowskip=4pt,
  columns=fullflexible,
}
\lstset{style=rstyle}

\newcommand{\status}[1]{\textcolor{dlnavy}{\textbf{[#1]}}}

\title{\vspace{-1.2cm}\color{dlnavy}\textbf{Disconnected Landlords --- Progress Pack}\\[2pt]
\large\normalcolor Call briefing, 10 July 2026 (LA geography, \texttt{Full-Sample})}
\author{Dmytro Kunchenko}
\date{Prepared 2026-07-10}

\begin{document}
\maketitle
\vspace{-0.6cm}

\noindent This note summarises the three active workstreams behind the 09 July
update --- the EPC data transition, the property-archetype reweighting method, and
the transaction-vintage gradient --- with the underlying formulae, the actual
estimating code, and the current exhibits. Two further tasks (attrition accounting;
foreign non-profits and tax havens) are appended for reference and are not the focus
of tomorrow's call. \textbf{Sign convention:} \texttt{current\_energy\_efficiency}
(SAP score) is higher~$=$~better; matched-pair effects are treated~$-$~control, so a
positive value means the ownership type is more efficient than its matched
private-individual control. \texttt{bad\_epc\_c} is the share of a group below EPC
band C (higher~$=$~worse).

\vspace{-0.1cm}
{\small\tableofcontents}
\clearpage

% ============================================================
\section{Data: moving to the new EPC release}
% ============================================================
\status{In progress}

The underlying EPC certificate dataset is being migrated to its newer release. The
schema is close to the old one, so most of the pipeline carries over unchanged. The
one complication is that ownership is observed from the OCOD/CCOD company-ownership
registers, which are a snapshot dated \textbf{January 2025}; certificates lodged
after that date could belong to a property that has since changed hands, and pairing
a post-snapshot EPC with a pre-snapshot owner risks misattributing the certificate to
the wrong landlord. The current, conservative fix is to cut the certificate feed off
at January 2025 as well, so both sources sit on the same clock. This only affects the
roughly 18-month tail; a cleaner long-run fix is to refresh EPC \emph{and} ownership
data together to a common, more recent vintage, which is worth planning for.

% ============================================================
\section{Property archetypes and reweighting}
% ============================================================
\status{Core method --- preliminary results}

\subsection{Motivation}
Different landlord types own systematically different property stock: offshore
owners concentrate in modern flats, councils in post-war semis. A raw treated-minus-
control gap therefore mixes two things --- a genuine \textbf{treatment effect} (how an
owner runs a property) and a \textbf{composition effect} (what kind of property that
owner happens to hold). The archetype exercise separates the two, so that a
like-for-like effect on a common reference stock can be reported alongside the raw
gap.

\subsection{Matched pairs}
The design is unchanged from the main regressions: 1:1 nearest-neighbour propensity
matching within exact strata. Categorical covariates --- property type, built form,
construction age band, main fuel --- are exact-matched, as is local authority;
habitable rooms and floor area are matched on nearest neighbour. Each stratum $s$
therefore contains exactly one treated and one control property.

\subsection{Step 1: recovering the pooled effect by differencing}
The baseline matching regression for property $i$ in subclass $s$ is
\begin{equation}
  Y_{is} = \tau\,\mathrm{Treat}_{is} + \gamma_1 \mathrm{Rooms}_{is}
           + \gamma_2 \mathrm{Floor}_{is} + \delta_s + \epsilon_{is}.
\end{equation}
Because each subclass holds one treated and one control unit, differencing the two
removes the subclass fixed effect $\delta_s$:
\begin{equation}
  \Delta Y_s = \tau + \gamma_1 \Delta \mathrm{Rooms}_s
             + \gamma_2 \Delta \mathrm{Floor}_s + \Delta\epsilon_s .
\end{equation}
This is estimated directly in \texttt{06c\_heterogeneous\_effects.R}:
\begin{lstlisting}
att <- feols(as.formula(paste0(dy, " ~ 1 + D_rooms + D_floor")),
             data = est, cluster = ~local_authority, lean = TRUE)
\end{lstlisting}
The intercept recovers the same pooled ATT as the headline PSM~+~subclass-FE
regression, without needing to build the subclass fixed effect at all.

\subsection{Step 2: cell-specific effects}
Splitting properties into archetype cells $c \in C$ (property type $\times$ built
form $\times$ age band $\times$ main fuel $\times$ floor-area tercile) and replacing
the single intercept with one per cell gives the conditional average treatment
effect (CATE) for each archetype:
\begin{equation}
  \Delta Y_s = \sum_{c\in C}\beta_c\,\mathbb{I}(\mathrm{Cell}_s=c)
             + \gamma_1 \Delta \mathrm{Rooms}_s
             + \gamma_2 \Delta \mathrm{Floor}_s + \Delta\epsilon_s .
\end{equation}
\begin{lstlisting}
est[, cell_f := factor(cell_est)]
cm <- feols(as.formula(paste0(dy, " ~ 0 + i(cell_f) + D_rooms + D_floor")),
            data = est, cluster = ~local_authority, lean = TRUE)
\end{lstlisting}

\subsection{Step 3: reweighting to a common stock}
Let $w_c$ be cell $c$'s share of the eligible private-rental control pool
($\sum_c w_c = 1$). The reweighted effect is the control-pool-weighted average of the
cell CATEs, with a clustered delta-method standard error:
\begin{equation}
  \tau_{rw} = \sum_{c\in C} w_c\,\beta_c ,
  \qquad
  \mathrm{SE}(\tau_{rw}) = \sqrt{W^{\!\top} V\, W},
\end{equation}
where $V$ is the local-authority-cluster-robust variance-covariance matrix of
$\hat\beta$ from the Step~2 regression.
\begin{lstlisting}
w_norm <- w_raw / cov_full
rw_row[, tau_rw_full := sum(w_norm * b[real_idx])]
rw_row[, tau_rw_full_se := sqrt(as.numeric(
  t(w_norm) %*% V[real_idx, real_idx, drop = FALSE] %*% w_norm))]
rw_row[, composition_effect := att_coef - tau_rw_full]
\end{lstlisting}

\subsection{Step 4: the composition effect}
\begin{equation}
  \underbrace{\tau}_{\text{pooled ATT}} - \underbrace{\tau_{rw}}_{\text{like-for-like}}
  \;=\; \text{composition effect}.
\end{equation}
This is the part of the headline gap explained purely by \emph{which} properties the
treated group owns, holding fixed how each individual archetype is treated.

\subsection{Worked example: two real archetypes}
Rather than list every cell, two contrasting archetypes make the mechanics concrete.
Both are drawn from the current \texttt{06c} output for the EPC-score outcome,
LA geography, baseline spec.

\vspace{2pt}
\noindent\textbf{The modal archetype (A1): Edwardian gas mid-terrace.}
The single most common cell in the control pool --- a mid-terrace house, built
1900--1929, on mains gas, medium floor area --- covers \textbf{4.04\% of the control
pool} (94,446 of 2,339,552 eligible properties) and has a control-pool baseline of
62.1 SAP points. Its cell CATEs: for-profit owners are essentially neutral here
($\beta = -0.37$), while non-profit owners are \textbf{4.38 points more efficient}
than the matched private-individual control on exactly this property type.

\vspace{2pt}
\noindent\textbf{An interesting minority archetype (A4): the period electric flat.}
A mid-terrace flat built 1900--1929 on electricity, small floor area, covers only
\textbf{0.76\% of the control pool} (17,887 properties) but has a much lower
baseline (56.2 SAP --- electric heating is expensive and inefficient in older stock)
and far larger CATEs: for-profit $+2.33$, non-profit $+7.12$, tax haven $+7.16$. This
is the closest \emph{estimable} analogue to the property type that offshore owners
actually concentrate in --- modern (2003--2006) electric flats. That true modal cell
for offshore ownership never clears the archetype whitelist (the top 100 cells by
control-pool frequency): no 2000s-built electric flat appears among them, only
pre-1950 ones. In other words, the stock offshore landlords typically hold is too
rare in the wider private-rental pool to get its own matched CATE at all --- which is
exactly why their composition effect is so large: the like-for-like comparison has
to fall back on the common stock they share with everyone else.

\vspace{2pt}
\noindent\textbf{Building the weighted estimate.}
$\tau_{rw}$ sums $w_c\beta_c$ over roughly 90--100 whitelisted cells; A1 and A4 are
two terms in that sum. For non-profit ownership, A1 contributes
$0.0404 \times 4.38 \approx 0.18$ points and A4 contributes
$0.0076 \times 7.12 \approx 0.05$ points --- together under a tenth of the full
reweighted effect ($\tau_{rw} = +3.29$, against a raw pooled ATT of $+3.15$): no
single archetype drives the result, which is the point of reweighting over the whole
distribution rather than reading off one or two cells.

\subsection{Reading the exhibits}
The archetype matrix (Table~\ref{tab:hte_archetypes}, Figure~\ref{fig:hte_matrix})
and the forest plot (Figure~\ref{fig:hte_forest}) report the \emph{same} cell CATEs
$\beta_c$ from Step~2, just in two formats: the matrix as a heatmap by archetype and
ownership type, the forest as points with confidence intervals. The forest adds one
extra piece of information, the pooled ATT $\tau$ from Step~1, marked with a red
dashed line, so the two views can be checked against each other cell-for-cell.

\subsection{Results}

\begin{table}[htbp]
\centering
\caption{The Ten Most Common Control-Group Property Archetypes}
\label{tab:hte_archetypes}
\scriptsize
\begin{tabular}{rlllllrr}
\toprule
 & Property type & Built form & Age band & Main fuel & Floor tercile & Share & Cum. \\
\midrule
1 & House & Mid-Terrace & England and Wales: 1900-1929 & mains gas (not community) & medium & 4.0\% & 4.0\% \\
2 & House & Mid-Terrace & England and Wales: 1900-1929 & mains gas (not community) & large & 3.7\% & 7.7\% \\
3 & House & Semi-Detached & England and Wales: 1930-1949 & mains gas (not community) & large & 2.0\% & 9.7\% \\
4 & Flat & Mid-Terrace & England and Wales: 1900-1929 & mains gas (not community) & small & 2.0\% & 11.7\% \\
5 & House & Mid-Terrace & England and Wales: before 1900 & mains gas (not community) & medium & 1.7\% & 13.4\% \\
6 & House & Mid-Terrace & England and Wales: before 1900 & mains gas (not community) & large & 1.7\% & 15.1\% \\
7 & House & Semi-Detached & England and Wales: 1930-1949 & mains gas (not community) & medium & 1.7\% & 16.8\% \\
8 & House & Semi-Detached & England and Wales: 1950-1966 & mains gas (not community) & medium & 1.5\% & 18.3\% \\
9 & House & Mid-Terrace & England and Wales: 1930-1949 & mains gas (not community) & medium & 1.2\% & 19.6\% \\
10 & Flat & Mid-Terrace & England and Wales: before 1900 & mains gas (not community) & small & 1.2\% & 20.8\% \\
\bottomrule
\end{tabular}
\begin{minipage}{0.95\textwidth}
\vspace{4pt}
\footnotesize
\textit{Notes:} Archetypes are the most frequent covariate cells
(property type $\times$ built form $\times$ construction age band $\times$ main fuel
$\times$ floor-area tercile) in the eligible private-rental control pool, England-wide.
Floor-area terciles cut at 64 and 86 sqm on the control pool.
Share = share of the control pool in the cell; Cum. = cumulative share.
Source: \texttt{hte\_archetype\_definitions\_LA.csv} (script 06c).
\end{minipage}
\end{table}


\begin{table}[htbp]
\centering
\caption{Composition vs.\ Treatment: Reweighted Effects, EPC Score}
\label{tab:hte_reweighted_current_energy_efficiency}
\small
\begin{tabular}{lrrrrr}
\toprule
Treatment & ATT (SAP pts) & $\tau_{rw}$ (full) & $\tau_{rw}$ (10 arch.) & Composition & Wald $p$ \\
\midrule
For-profit (all) & +0.45*** & +0.16 (0.050) & -0.03 (0.064) & +0.29 & $<$0.001 \\
UK for-profit & +0.44*** & +0.15 (0.051) & -0.03 (0.065) & +0.29 & $<$0.001 \\
Foreign for-profit & +1.53*** & +0.64 (0.138) & +0.10 (0.167) & +0.88 & $<$0.001 \\
Non-profit & +3.15*** & +3.29 (0.119) & +3.58 (0.127) & -0.14 & $<$0.001 \\
Tax haven (all) & +1.50*** & +0.65 (0.141) & -0.04 (0.203) & +0.85 & $<$0.001 \\
Public sector & +3.38*** & +2.62 (0.158) & +2.73 (0.195) & +0.76 & $<$0.001 \\
\bottomrule
\end{tabular}
\begin{minipage}{0.95\textwidth}
\vspace{4pt}
\footnotesize
\textit{Notes:} Matched-pair design, Baseline spec, base core, LA-clustered.
ATT = pooled pair-difference estimate (identical to the PSM + Subclass FE
coefficient). $\tau_{rw}$ = cell CATEs reweighted to the control-pool covariate
distribution (delta-method SEs in parentheses); full = all estimated cells,
10 arch. = the ten headline archetypes. Composition = ATT $-$ $\tau_{rw}$ (full):
the part of the pooled effect attributable to the treated group's covariate mix.
Wald $p$: test that all cell CATEs are equal.

Source: \texttt{hte\_reweighted\_LA.csv} (script 06c).
\end{minipage}
\end{table}


Public sector and non-profit owners are genuinely more efficient than private
rental like-for-like: their reweighted effect is as large as, or larger than, the
raw ATT. For offshore and haven owners, most of the raw efficiency advantage is
composition --- reweighted to a common stock, the premium collapses toward zero.

\begin{table}[htbp]
\centering
\caption{Typical (Modal) Property by Ownership Type}
\label{tab:hte_typical_property}
\scriptsize
\begin{tabular}{lllllrrc}
\toprule
Ownership & Property type & Built form & Age band & Main fuel & Share & Mean EPC & Estimable \\
\midrule
For-Profit & House & Mid-Terrace & 1900-1929 & Mains gas & 2.2\% & 67.1 & yes \\
UK For-Profit & House & Mid-Terrace & 1900-1929 & Mains gas & 2.3\% & 66.9 & yes \\
Foreign For-Profit & Flat & Mid-Terrace & 2003-2006 & Electricity & 3.0\% & 71.6 & no \\
Non-Profit & House & Semi-Detached & 1950-1966 & Mains gas & 2.9\% & 69.9 & yes \\
Public Sector & House & Semi-Detached & 1950-1966 & Mains gas & 3.7\% & 68.7 & yes \\
Tax Haven & Flat & Mid-Terrace & 2003-2006 & Electricity & 2.9\% & 71.6 & no \\
Tax Haven For-Profit & Flat & Mid-Terrace & 2003-2006 & Electricity & 3.0\% & 71.8 & no \\
Tax Haven Non-Profit & Flat & Mid-Terrace & 1900-1929 & Mains gas & 4.9\% & 67.5 & yes \\
Caribbean Haven & Flat & Mid-Terrace & 2003-2006 & Electricity & 2.2\% & 69.1 & no \\
British Haven & Flat & Mid-Terrace & 2003-2006 & Electricity & 2.9\% & 71.9 & no \\
European Haven & Flat & Mid-Terrace & 2003-2006 & Electricity & 3.0\% & 72.3 & no \\
Other Haven & Flat & Mid-Terrace & 2003-2006 & Electricity & 5.1\% & 68.0 & no \\
Foreign Non-Profit & Flat & Mid-Terrace & 1930-1949 & Mains gas & 7.9\% & 67.7 & no \\
UK Non-Profit & House & Semi-Detached & 1950-1966 & Mains gas & 2.9\% & 69.9 & yes \\
\bottomrule
\end{tabular}
\begin{minipage}{0.95\textwidth}
\vspace{4pt}
\footnotesize
\textit{Notes:} For each ownership type, the single most common covariate cell among
its \emph{treated} eligible properties (property type $\times$ built form $\times$ age
band $\times$ main fuel $\times$ floor tercile; tercile omitted here for space).
Share = share of the treated group falling in that modal cell. Mean EPC = mean current
energy-efficiency score of the treated group. Estimable = the modal cell is in the
archetype whitelist \emph{and} has $\geq 50$ matched pairs for that ownership type, so a
matched CATE exists for it; distinctive treated cells (e.g.\ prime flats) often lack
matched support and so are not estimable, which is why the matched archetypes reflect
the common control stock instead. Source: \texttt{hte\_typical\_property\_LA.csv} (06c).
\end{minipage}
\end{table}


\begin{figure}[H]
  \centering
  \includegraphics[width=\textwidth,height=0.78\textheight,keepaspectratio]%
    {output/figures/hte_archetype_forest.pdf}
  \caption{Archetype CATEs ($\beta_c$) for six ownership types and two outcomes; red
  dashed line = pooled ATT $\tau$. Rows are the archetypes of
  Table~\ref{tab:hte_archetypes}.}
  \label{fig:hte_forest}
\end{figure}

\begin{figure}[H]
  \centering
  \includegraphics[width=\textwidth,height=0.72\textheight,keepaspectratio]%
    {output/figures/hte_archetype_matrix_LA.pdf}
  \caption{Archetype~$\times$~ownership CATEs as a heatmap (green = more efficient,
  purple = less). Baseline EPC in each row label; dashed outline = suspect support;
  grey = no matched support. Council-stock references shown separately.}
  \label{fig:hte_matrix}
\end{figure}

% ============================================================
\section{Transaction vintage and retrofitting}
% ============================================================
\status{Preliminary --- descriptive}

The same pair-difference machinery, applied to time since a property last
transacted, asks whether the ownership efficiency gap fades with time. Pairs are
additionally exact-matched on vintage (both lodgement year and sale year), so time-
since-sale is constant within a pair; sales postdating the EPC are dropped.

\subsection{The estimator ladder}
Three estimators of the same penalty are compared, from most to least confounded by
composition. \emph{Raw}: an unmatched per-bin gap that still reflects whichever
properties happen to transact in that window,
\begin{lstlisting}
feols(oc ~ i(ts_bin_f, treat_var) + ts_bin_f,
      data = est, cluster = ~local_authority, lean = TRUE)
\end{lstlisting}
\emph{Adjusted}: the same regression with property-type, built-form, age-band and
main-fuel fixed effects plus floor area and habitable rooms as controls --- still
unmatched, but composition-adjusted,
\begin{lstlisting}
feols(oc ~ i(ts_bin_f, treat_var) + ts_bin_f + total_floor_area + number_habitable_rooms
      | property_type + built_form + construction_age_band + main_fuel,
      data = est, cluster = ~local_authority, lean = TRUE)
\end{lstlisting}
\emph{Matched}: the pair-difference gradient, which strips composition out entirely,
\begin{lstlisting}
feols(as.formula(paste0(dy, " ~ 0 + i(ts_bin_f)")),
      data = est, cluster = ~local_authority, lean = TRUE)
\end{lstlisting}
Raw $\to$ adjusted $\to$ matched typically shrinks by an order of magnitude: for
foreign for-profit and tax-haven owners the unmatched gap is around 10 SAP points,
falls to around 2.5 once composition-adjusted, and to under 1 once matched --- the
raw comparison is dominated by the fact that offshore owners hold newer, better
stock, exactly as in the archetype exercise above.

\begin{table}[htbp]
\centering
\caption{EPC-Score Penalty by Time Since Last Transaction}
\label{tab:txn_vintage}
\small
\begin{tabular}{lrrrrr}
\toprule
Years since sale & For-profit & Foreign FP & Tax haven & TH for-profit & Non-profit \\
\midrule
0-1 & -0.40*** (0.08) & +0.67 (0.49) & +0.57 (0.49) & +0.65 (0.51) & +1.15*** (0.29) \\
2-4 & +0.40*** (0.08) & +0.77 (0.48) & +0.97* (0.46) & +0.98* (0.46) & +1.47*** (0.33) \\
5-9 & +0.25** (0.09) & -0.01 (0.31) & +0.15 (0.33) & +0.11 (0.33) & +2.34*** (0.26) \\
10-14 & +0.31*** (0.06) & +0.42 (0.27) & +0.39 (0.27) & +0.36 (0.27) & +2.50*** (0.21) \\
15+ & +0.43*** (0.07) & +0.41 (0.34) & +0.53 (0.33) & +0.49 (0.34) & +1.85*** (0.15) \\
\bottomrule
\end{tabular}
\begin{minipage}{0.95\textwidth}
\vspace{4pt}
\footnotesize
\textit{Notes:} Mean matched-pair difference in EPC score (treated $-$ control)
by years between the property's most recent sale and its EPC lodgement.
Price Paid specification, ppd matching core: pairs exact-match on both
\texttt{lodgement\_year} and \texttt{ppd\_year\_transfer}, so the vintage is
constant within pair. LA-clustered SEs in parentheses.
$^{***}p<0.001$; $^{**}p<0.01$; $^{*}p<0.05$.
\textbf{Descriptive:} selection into transaction timing is not random; this is a
gradient shape, not identification of retrofit dynamics. Sales postdating the EPC
are excluded (share reported in \texttt{txn\_vintage\_gradient\_LA.csv}).
Source: script 06d.
\end{minipage}
\end{table}


\begin{table}[htbp]
\centering
\caption{Estimator Ladder: EPC-Score Ownership Penalty by Transaction Vintage}
\label{tab:txn_vintage_simple}
\small
\begin{tabular}{llrrr}
\toprule
Ownership & Years since sale & Simple (raw) & Simple (adj.) & Matched \\
\midrule
For-profit & 0-1 & +4.06 & +0.66 & -0.40 \\
 & 2-4 & +5.44 & +1.26 & +0.40 \\
 & 5-9 & +5.45 & +1.09 & +0.25 \\
 & 10-14 & +5.15 & +0.63 & +0.31 \\
 & 15+ & +3.57 & +0.65 & +0.43 \\
\midrule
Foreign for-profit & 0-1 & +9.90 & +3.02 & +0.67 \\
 & 2-4 & +9.96 & +2.64 & +0.77 \\
 & 5-9 & +10.17 & +2.56 & -0.01 \\
 & 10-14 & +9.97 & +2.43 & +0.42 \\
 & 15+ & +6.58 & +2.40 & +0.41 \\
\midrule
Tax haven & 0-1 & +9.94 & +3.00 & +0.57 \\
 & 2-4 & +9.95 & +2.62 & +0.97 \\
 & 5-9 & +10.26 & +2.60 & +0.15 \\
 & 10-14 & +9.95 & +2.39 & +0.39 \\
 & 15+ & +6.66 & +2.46 & +0.53 \\
\midrule
Tax-haven for-profit & 0-1 & +10.01 & +3.03 & +0.65 \\
 & 2-4 & +10.01 & +2.63 & +0.98 \\
 & 5-9 & +10.29 & +2.62 & +0.11 \\
 & 10-14 & +10.02 & +2.42 & +0.36 \\
 & 15+ & +6.74 & +2.47 & +0.49 \\
\midrule
Non-profit & 0-1 & +6.88 & +2.10 & +1.15 \\
 & 2-4 & +6.79 & +2.19 & +1.47 \\
 & 5-9 & +6.16 & +2.18 & +2.34 \\
 & 10-14 & +5.16 & +1.97 & +2.50 \\
 & 15+ & +4.05 & +1.81 & +1.85 \\
\bottomrule
\end{tabular}
\begin{minipage}{0.95\textwidth}
\vspace{4pt}
\footnotesize
\textit{Notes:} Ownership penalty in EPC score (treated $-$ control) by years
between the property's most recent sale and its EPC lodgement. \emph{Simple (raw)}:
unmatched per-bin gap from $y \sim \mathrm{i}(\text{bin},\text{owner}) + \text{bin}$.
\emph{Simple (adjusted)}: same, plus property-type, built-form, age-band, main-fuel
fixed effects and floor area / habitable rooms (unmatched but composition-adjusted).
\emph{Matched}: mean pair difference (Price Paid spec, ppd core; exact-matched on
vintage). Raw includes composition by vintage; matched strips it; adjusted sits
between. LA-clustered SEs. \textbf{Descriptive:} selection into transaction timing is
not random. Source: script 06d.
\end{minipage}
\end{table}


For-profit and offshore owners show a small, roughly flat matched gap across
vintage bins --- once composition is stripped out there is little vintage story left
to tell. Non-profit owners show a larger matched gap that peaks at 5--14 years since
sale, consistent with a retrofit cycle on retained stock. Selection into transaction
timing is not random, so this is a descriptive gradient shape, not an identified
retrofit dynamic.

\begin{figure}[H]
  \centering
  \includegraphics[width=\textwidth,height=0.6\textheight,keepaspectratio]%
    {output/figures/txn_vintage_penalty_LA.pdf}
  \caption{Matched-pair ownership penalty by years since last sale (levels; 95\% CI
  ribbons; panel per outcome).}
  \label{fig:txn_penalty}
\end{figure}

\begin{figure}[H]
  \centering
  \includegraphics[width=\textwidth,height=0.6\textheight,keepaspectratio]%
    {output/figures/txn_vintage_simple_LA.pdf}
  \caption{Estimator ladder: unmatched raw and covariate-adjusted vs.\ matched
  pair-difference penalty (free $y$-scale per panel).}
  \label{fig:txn_ladder}
\end{figure}

% ============================================================
\section{Matching attrition \textnormal{\small(reference)}}
% ============================================================
\status{Complete --- not the focus of this call}

Property- and local-authority-level survival through the matching funnel (eligible
$\to$ complete covariate cases $\to$ size gate $\to$ 1:1 propensity match $\to$
caliper filters), computed in \texttt{05b\_compute\_balance.R}. Match rates run
27--45\%; the caliper is aggressive for the large for-profit pool but retains most
matches for the smaller, more distinctive offshore pools. Local-authority coverage
stays close to complete throughout.

\begin{table}[htbp]
\centering
\caption{Matching Attrition Funnel: Treated Units and LA Survival by Stage}
\label{tab:attrition_funnel}
\small
\begin{tabular}{lrrrrr}
\toprule
Stage & For-profit & Foreign FP & Tax haven & Non-profit & Public \\
\midrule
\multicolumn{6}{l}{\textit{Panel A: Baseline spec, base core}} \\
Eligible (raw, core sample) & 3,027,035 & 174,526 & 174,947 & 2,489,856 & 1,435,432 \\
Complete covariate cases & 2,282,352 & 90,252 & 90,336 & 2,002,000 & 1,361,592 \\
Passed size gate & 2,282,352 & 90,109 & 90,202 & 2,002,000 & 1,361,460 \\
Matched (1:1 PSM) & 913,094 & 40,233 & 40,385 & 559,071 & 360,774 \\
Caliper PS $\leq$ 0.2 & 23,356 & 34,412 & 34,314 & 37,101 & 26,631 \\
Caliper PS $\leq$ 0.1 & 4,214 & 25,625 & 25,503 & 6,728 & 8,638 \\
\addlinespace
LAs at start & 346 & 346 & 346 & 346 & 346 \\
LAs with eligible sample & 346 & 345 & 346 & 346 & 346 \\
LAs passing gate & 346 & 319 & 319 & 346 & 325 \\
LAs with matches & 346 & 325 & 324 & 346 & 331 \\
LAs in headline regression & 346 & 329 & 326 & 346 & 338 \\
\midrule
\multicolumn{6}{l}{\textit{Panel B: Price Paid spec, ppd core}} \\
Eligible (raw, core sample) & 2,000,667 & 119,985 & 120,657 & 332,225 & 278,108 \\
Complete covariate cases & 1,638,777 & 65,014 & 65,245 & 261,186 & 261,863 \\
Passed size gate & 1,637,430 & 64,887 & 65,113 & 261,135 & 260,469 \\
Matched (1:1 PSM) & 218,576 & 5,423 & 5,290 & 27,607 & 21,243 \\
Caliper PS $\leq$ 0.2 & 13,529 & 4,033 & 3,873 & 16,852 & 8,393 \\
Caliper PS $\leq$ 0.1 & 3,959 & 3,012 & 2,859 & 9,219 & 4,683 \\
\addlinespace
LAs at start & 346 & 346 & 346 & 346 & 346 \\
LAs with eligible sample & 346 & 343 & 344 & 345 & 337 \\
LAs passing gate & 344 & 320 & 317 & 344 & 295 \\
LAs with matches & 345 & 262 & 255 & 343 & 275 \\
LAs in headline regression & 345 & 269 & 261 & 344 & 278 \\
\bottomrule
\end{tabular}
\begin{minipage}{0.95\textwidth}
\vspace{4pt}
\footnotesize
\textit{Notes:} Counts of treated properties surviving each stage of the matching
funnel, pooled across local authorities, plus the number of LAs surviving each stage.
Eligible = non-missing treatment indicator within the core sample. Complete cases =
non-missing matching covariates. Size gate = per-LA minimum treated/control counts
(sparse cores: $>5$ treated, $>25$ control; other cores: $>10$ treated, $>50$ control).
Matched = 1:1 nearest-neighbour propensity-score matches with exact strata.
Calipers are post-hoc propensity-score filters as in the PS$\leq$0.2 / PS$\leq$0.1 models.
``LAs in headline regression'' = distinct-LA cluster count of the PSM + Subclass FE model
(outcome \texttt{current\_energy\_efficiency}).
Source: \texttt{attrition\_funnel\_LA.csv} (script 05b).
\end{minipage}
\end{table}


% ============================================================
\section{Foreign non-profits and tax havens \textnormal{\small(reference)}}
% ============================================================
\status{Descriptive --- not the focus of this call}

A descriptive profile of offshore and foreign non-profit ownership from company-level
public register data (OCOD/CCOD plus jurisdiction). The tax-haven universe is
dominated by Crown-Dependency (Guernsey/Jersey) ground-rent vehicles --- large
freehold-reversion portfolios spanning over a hundred local authorities with good
EPCs --- rather than the disinvesting offshore landlords the hypothesis anticipates.
A smaller, genuinely opaque tail (British Virgin Islands, Gibraltar, and the wider
Caribbean/other haven classes) holds worse and more expensive prime stock. Full
tables are in \texttt{task\_exhibits.pdf}; named case studies are drafted in
\texttt{task\_4\_offshore\_nonprofit\_profiles.tex}.

% ============================================================
\section{For discussion tomorrow}
% ============================================================
\begin{enumerate}
  \item \textbf{Data:} accept the January-2025 cutoff, or commit to refreshing EPC and
        ownership data to a common recent vintage?
  \item \textbf{Reweighting:} is the control pool the right reference distribution for
        $w_c$, or should it be the treated pool, or a national stock benchmark?
  \item \textbf{Scope:} the pair-difference collapse is cheap and exact --- worth
        applying to the main regressions more broadly?
  \item \textbf{Vintage:} keep the gradient descriptive, or is it worth pursuing an
        identified retrofit design?
\end{enumerate}

% ============================================================
\appendix
\section{Notation and code pointers}
\begin{itemize}
  \item $i$ property, $s$ matched pair/subclass, $c$ archetype cell.
  \item $\tau$ pooled ATT; $\beta_c$ cell CATE; $w_c$ control-pool share;
        $\tau_{rw} = W^\top\beta$.
  \item $V$ local-authority-cluster-robust VCOV of $\hat\beta$;
        $\mathrm{SE}(\tau_{rw}) = \sqrt{W^\top V W}$.
  \item Archetype = property type $\times$ built form $\times$ age band $\times$ main
        fuel $\times$ floor tercile (no geography; both units of a pair share an LA
        by construction).
\end{itemize}
\textbf{Scripts:} \texttt{05b\_compute\_balance.R} (attrition),
\texttt{06c\_heterogeneous\_effects.R} (archetypes and reweighting),
\texttt{06d\_transaction\_vintage.R} (vintage ladder),
\texttt{archetype\_names.R} (narrative labels). Method note:
\texttt{docs/06c\_reweighting\_methodology.md}.

\end{document}
